% Sección P (añadida en la versión revisada del 17-11-2026). Lleva la letra P, como en el cuaderno, para que las
% secciones siguientes conserven su número: otros documentos de la entrega las citan por número.
% Todas las cifras son macros de generado/corto_plazo.tex, que informe/cifras_corto_plazo.py lee de
% pronostico_corto_plazo/panel/ y pronostico_corto_plazo/resultados/ (ninguna está escrita a mano). Los horizontes van
% en romanos en el nombre de la macro: XXX = 30, LX = 60 y CXX = 120 minutos.
\renewcommand{\thesection}{P}\renewcommand{\theHsection}{P}
\section{El problema del enunciado: pronóstico a 30, 60 y 120 minutos por ruta, sentido y segmento}
\label{sec:corto-plazo}
\renewcommand{\thesection}{\arabic{section}}\renewcommand{\theHsection}{\arabic{section}}
\addtocounter{section}{-1}
% Sus tablas y figuras se numeran P.1, P.2...: las del resto del informe conservan también su número.
\newcounter{cptabla}\newcounter{cpfigura}
\setcounter{cptabla}{\value{table}}\setcounter{cpfigura}{\value{figure}}
\setcounter{table}{0}\setcounter{figure}{0}
\renewcommand{\thetable}{P.\arabic{table}}\renewcommand{\theHtable}{P.\arabic{table}}
\renewcommand{\thefigure}{P.\arabic{figure}}\renewcommand{\theHfigure}{P.\arabic{figure}}

La especificación del reto pide \comi{predecir la demanda para horizontes configurables de 30, 60 y 120 minutos}
\comi{por ruta, sentido, segmento y franja horaria}, con \comi{división cronológica en entrenamiento, validación y
prueba reservada}. La primera versión de este informe resolvía otro problema: el pronóstico a una semana por ruta y
hora de salida, que es el que necesita la programación semanal (secciones~\ref{sec:metodologia} y
\ref{sec:resultados}). La retroalimentación del curso señaló tres faltas: ese no era el problema pedido; los modelos
se habían elegido y reportado sobre los mismos meses de 2025, y el modelo por segmento se había entrenado con 37 días.
Esta sección resuelve el problema pedido y corrige las tres. El código, el protocolo y los resultados están en
\texttt{pronostico\_corto\_plazo/}.

\paragraph{Datos.} La tabla \texttt{timbradas\_segmentos} de la base solo cubre 58 días de 2023; de ahí los 37 días
de entrenamiento. Para tener los tres años exportamos, en solo lectura, los \CPEventos{} eventos de conteo con subidas
y asignamos cada uno al segmento oficial más cercano de su ruta, que es la regla con que Extreme llena esa tabla. En
los 58 días comunes la regla propia da las mismas \CPVerSubidas{} subidas, con el \CPVerMismoSegmento~\% en el mismo
segmento y el \CPVerMismoTotal~\% de los viajes con el mismo total. El panel tiene \CPSubidas{} subidas por ruta, día,
bloque de 15 minutos y segmento, en los 386 segmentos de 500~m de las tres rutas. Un día de una ruta con menos del
5~\% de las subidas habituales se trata como día sin datos: son los once últimos días de mes de 2025 que la base no
trae. El sentido no viene en los datos: \comi{ida} es la primera mitad del ciclo y \comi{regreso} la segunda. La
regla del segmento más cercano no mira en qué parte del viaje va el bus y, donde la ida y el regreso comparten calle,
deja una parte de las subidas en el sentido contrario: cerca del \CPSentOtroSil~\% en SILENCIO, del \CPSentOtroCra~\%
en CRA~54 y del \CPSentOtroCl~\% en CL~72, medido con el tiempo de viaje que lleva el bus en cada subida.

\paragraph{Problema y protocolo.} Cada media hora, de 04:00 a 21:00, se pronostican las subidas de los 30, 60 y 120
minutos siguientes con lo ocurrido hasta ese instante, en tres niveles: la ruta, el sentido y el segmento. Ninguna
variable usa nada posterior: una comprobación automática cambia todo lo que ocurre después de un instante y verifica
que las variables y las líneas base anteriores no se mueven (\CPFugaComprobaciones{} comprobaciones,
\CPFugaDiferencias{} diferencias). El
protocolo se escribió y se confirmó en el repositorio antes de correr la validación: el historial de la entrega trae, en este
orden, el commit del protocolo (\texttt{9844dae}), el de la validación con lo elegido (\texttt{fa04afb}) y el de la
prueba (\texttt{546cffa}), los tres del 17 de noviembre de 2026. Se entrena con 2023 y 2024; todo se elige
con enero a junio de 2025 (variables, hiperparámetros, estructura y banda), y la prueba, de julio a diciembre de 2025,
se calcula una sola vez con lo elegido, reentrenado hasta junio. Las líneas base son las del enunciado, el promedio
por día y hora (la misma ventana del mismo día de la semana en las ocho semanas anteriores) y la última semana, más la
persistencia (lo que subió en los mismos minutos justo antes). El candidato es un LightGBM con pérdida de Poisson por
nivel, ruta y horizonte. Los meses de 2025 ya se conocían por los experimentos anteriores, con otro objetivo, así que
la prueba no es ciega en sentido estricto; ninguna decisión de este experimento usa julio a diciembre.

\paragraph{Lo elegido.} En la validación y por ruta, el WAPE medio de los tres horizontes es \CPValRutaProm~\% con el
promedio por día y hora; \CPValRutaSemanal~\% con el modelo que solo usa el calendario y las semanas anteriores, que es
lo que se sabe con siete días de antelación; \CPValRutaConHoy~\% al añadir lo observado en el día, y
\CPValRutaCompleto~\% al añadir las salidas ya despachadas. Se eligió \CPEligeConjuntos{} y, entre tres
configuraciones de árboles fijadas de antemano, \CPEligeHiper. Por sentido se eligió \CPEligeSentEstructura{}
(\CPValSentDirecto~\% con el modelo directo y \CPValSentSuma~\% sumando segmentos) y por segmento,
\CPEligeSegEstructura{} (\CPValSegDirecto~\% directo, \CPValSegConciliado~\% conciliado con el pronóstico de la ruta y
\CPValSegReparto~\% repartiendo el pronóstico de la ruta con las fracciones históricas).

\begin{table}[t]
  \centering\footnotesize
  \setlength{\tabcolsep}{3.5pt}
  \caption{Prueba reservada (julio a diciembre de 2025): WAPE en \% de cada pronóstico sobre las mismas celdas.
  \comi{Promedio} es el promedio por día y hora. La diferencia es la del modelo con ese promedio, en puntos, con IC
  del 95~\% por \emph{bootstrap} de bloques de siete días. La banda P10--P90 es conformal y su cobertura objetivo es
  80~\%.}
  \label{tab:cp-prueba}
  \begin{tabular}{@{}l r r r r r l r@{}}
    \toprule
    Nivel & Horizonte & Promedio & Últ.\ semana & Persistencia & \textbf{Modelo} & Diferencia [IC 95~\%] & Cobertura \\
    \midrule
    Ruta & 30 min & \CPRutaPromXXX & \CPRutaSemXXX & \CPRutaPersXXX & \textbf{\CPRutaModXXX} & \CPRutaDifXXX{} [\CPRutaDifICXXX] & \CPRutaCobXXX \\
    Ruta & 60 min & \CPRutaPromLX & \CPRutaSemLX & \CPRutaPersLX & \textbf{\CPRutaModLX} & \CPRutaDifLX{} [\CPRutaDifICLX] & \CPRutaCobLX \\
    Ruta & 120 min & \CPRutaPromCXX & \CPRutaSemCXX & \CPRutaPersCXX & \textbf{\CPRutaModCXX} & \CPRutaDifCXX{} [\CPRutaDifICCXX] & \CPRutaCobCXX \\
    \midrule
    Sentido & 30 min & \CPSentPromXXX & \CPSentSemXXX & \CPSentPersXXX & \textbf{\CPSentModXXX} & \CPSentDifXXX{} [\CPSentDifICXXX] & \CPSentCobXXX \\
    Sentido & 60 min & \CPSentPromLX & \CPSentSemLX & \CPSentPersLX & \textbf{\CPSentModLX} & \CPSentDifLX{} [\CPSentDifICLX] & \CPSentCobLX \\
    Sentido & 120 min & \CPSentPromCXX & \CPSentSemCXX & \CPSentPersCXX & \textbf{\CPSentModCXX} & \CPSentDifCXX{} [\CPSentDifICCXX] & \CPSentCobCXX \\
    \midrule
    Segmento & 30 min & \CPSegPromXXX & \CPSegSemXXX & \CPSegPersXXX & \textbf{\CPSegModXXX} & \CPSegDifXXX{} [\CPSegDifICXXX] & \CPSegCobXXX \\
    Segmento & 60 min & \CPSegPromLX & \CPSegSemLX & \CPSegPersLX & \textbf{\CPSegModLX} & \CPSegDifLX{} [\CPSegDifICLX] & \CPSegCobLX \\
    Segmento & 120 min & \CPSegPromCXX & \CPSegSemCXX & \CPSegPersCXX & \textbf{\CPSegModCXX} & \CPSegDifCXX{} [\CPSegDifICCXX] & \CPSegCobCXX \\
    \bottomrule
  \end{tabular}
  \fuente{\texttt{pronostico\_corto\_plazo/resultados/prueba.json}. \CPDiasPrueba{} días, de los que
  \CPDiasPruebaConDatos{} tienen datos en cada ruta; 35 instantes por día.}
\end{table}

\paragraph{Resultados.} En la prueba, por ruta y a 60 minutos, el modelo se equivoca en \CPRutaModLX~\% de las
subidas y el promedio por día y hora en \CPRutaPromLX~\% (tabla~\ref{tab:cp-prueba}): quita el \CPRutaQuitaLX~\% del
error, y el intervalo de la diferencia \CPRutaDifFraseLX. A 30 minutos el modelo da \CPRutaModXXX~\% y a 120,
\CPRutaModCXX~\%. Por ruta, a 60 minutos, el error es \CPCraModLX~\% en CRA~54, \CPSilModLX~\% en SILENCIO y
\CPClModLX~\% en CL~72 (el promedio: \CPCraPromLX, \CPSilPromLX{} y \CPClPromLX~\%). En las horas pico el modelo
falla \CPRutaPicoModLX~\% y en el valle \CPRutaValleModLX~\%, frente a \CPRutaPicoPromLX{} y \CPRutaVallePromLX~\% del
promedio. En días laborales da \CPRutaLabModLX~\%, en sábados \CPRutaSabModLX~\% y en domingos y festivos
\CPRutaDomModLX~\% (el promedio: \CPRutaLabPromLX, \CPRutaSabPromLX{} y \CPRutaDomPromLX~\%). Ese último grupo mezcla dos
casos, que un análisis posterior separa: en los \CPFestDias{} festivos el promedio pronostica un día corriente y
falla \CPFestPromLX~\%, frente a \CPFestModLX~\% del modelo; en los \CPDomSoloDias{} domingos corrientes el modelo
\emph{no} mejora al promedio (\CPDomSoloModLX{} frente a \CPDomSoloPromLX~\%; a 120 minutos, \CPDomSoloModCXX{}
frente a \CPDomSoloPromCXX~\%), y la pérdida está en CRA~54 (\CPDomSoloCraModLX{} frente a \CPDomSoloCraPromLX~\%).
El MAE a 60 minutos es de \CPRutaMAELX{} subidas sobre una media de \CPRutaMediaLX, el RMSE de \CPRutaRMSELX{} y el sMAPE de
\CPRutaSMAPELX~\%; el sesgo es de \CPRutaSesgoLX~\%. Mes a mes, el error a 60 minutos va de \CPRutaMesMinLX{} a
\CPRutaMesMaxLX~\%.

\begin{figure}[t]
  \centering
  \includegraphics[width=0.98\linewidth]{fig_cp_segmentos}
  \caption{Subidas por ventana de 60 minutos en cada segmento de 500~m a lo largo del recorrido, en la prueba: lo que
  subió (barras) y lo pronosticado (línea). La línea punteada separa la ida del regreso. En CRA~54, la barra sin línea
  cerca del km~2 es el registro imposible del \CPAnomFecha.}
  \label{fig:cp-segmentos}
  \fuente{\texttt{pronostico\_corto\_plazo/resultados/segmentos\_prueba.csv} (\texttt{figuras.py}).}
\end{figure}

\paragraph{Lo espacial.} El modelo por segmento se entrenó con todo lo anterior a julio de 2025 y se probó en
\CPSegFilasLX{} celdas de segmento e instante que no vio. A 60 minutos se equivoca en \CPSegModLX~\% y el promedio por
día y hora en \CPSegPromLX~\%; el error es alto en los dos porque cada celda trae pocas subidas, \CPSegMediaLX{} en
promedio. El modelo gana al promedio en \CPSegGanaLX{} de los \CPSegConSubidasLX{} segmentos con subidas, que llevan
el \CPSegGanaSubidasLX~\% de las subidas. El valor absoluto del sesgo por segmento tiene una mediana de
\CPSegSesgoMedianoLX~\% y, ponderado por subidas, una media de \CPSegSesgoPonderadoLX~\%
(figura~\ref{fig:cp-segmentos}). Por sentido el modelo da \CPSentModLX~\% frente a \CPSentPromLX~\% del promedio.

\begin{figure}[t]
  \centering
  \includegraphics[width=0.72\linewidth]{fig_cp_informacion}
  \caption{WAPE por ruta en la prueba según la información que usa el pronóstico: el promedio por día y hora, el
  modelo con lo que se sabe siete días antes y el modelo elegido, sin retraso y con 15 y 30 minutos de retraso en la
  captura.}
  \label{fig:cp-informacion}
  \fuente{\texttt{pronostico\_corto\_plazo/resultados/prueba.json}, clave \texttt{secundarios}.}
\end{figure}

\paragraph{Lo observado en el día.} El protocolo declaró de antemano dos análisis por ruta. El mismo modelo con solo
lo que se sabe siete días antes se equivoca en \CPSemanalLX~\% a 60 minutos; con lo observado hasta el instante, en
\CPRutaModLX~\% (IC del 95~\% de la diferencia: \CPSemanalDifICLX{} puntos). Esa diferencia es lo que le faltaba al
pronóstico a una semana para ser el del enunciado. Con 15 y 30 minutos de retraso en la captura, el error a 60
minutos es de \CPRetrasoQuinceLX{} y \CPRetrasoTreintaLX~\% (figura~\ref{fig:cp-informacion}).

\paragraph{Un registro imposible.} El RMSE por ruta es varias veces mayor que el MAE por un registro de conteo
imposible que está en la base: el \CPAnomFecha{} a las \CPAnomHora, \CPAnomSubidas{} subidas en un solo segmento de
\CPAnomRuta{} en 15 minutos (lo habitual en toda la ruta en ese bloque son \CPAnomMediana). Hay otro, de
\CPAnomOtroSubidas{} subidas, en el entrenamiento (\CPAnomOtroFecha). El protocolo no previó filtrarlos, así que las
cifras de la tabla~\ref{tab:cp-prueba} lo incluyen y no se cambian. En un análisis posterior, con las predicciones
guardadas y sin los \CPAnomDias{} días de \CPAnomRuta{} en que ese registro entra en el objetivo o en el promedio de
ocho semanas, el modelo se equivoca por ruta a 60 minutos en \CPSinRutaModLX~\% y el promedio en \CPSinRutaPromLX~\%
(diferencia de \CPSinRutaDifLX{} puntos, con IC del 95~\% de \CPSinRutaDifICLX), y el RMSE baja de \CPRutaRMSELX{} a
\CPSinRutaRMSELX. La conclusión no cambia.

\paragraph{Lo que falta.} Las subidas son demanda observada: dependen de los buses que pasaron. El modelo supone que
el conteo llega en vivo, y la latencia real de la captura no se conoce. Los dos registros imposibles siguen en los
datos del entrenamiento: al servir se excluyen. La API lo sirve por omisión, reproduce las \CPReproFilas{}
predicciones de esta prueba y lo alimenta con la lectura incremental.

% Fin de la sección P: las tablas y figuras siguientes retoman la numeración del informe.
\setcounter{table}{\value{cptabla}}\setcounter{figure}{\value{cpfigura}}
\renewcommand{\thetable}{\arabic{table}}\renewcommand{\theHtable}{\arabic{table}}
\renewcommand{\thefigure}{\arabic{figure}}\renewcommand{\theHfigure}{\arabic{figure}}
